% Options for packages loaded elsewhere
\PassOptionsToPackage{unicode}{hyperref}
\PassOptionsToPackage{hyphens}{url}
\documentclass[
]{article}
\usepackage{xcolor}
\usepackage{amsmath,amssymb}
\setcounter{secnumdepth}{-\maxdimen} % remove section numbering
\usepackage{iftex}
\ifPDFTeX
  \usepackage[T1]{fontenc}
  \usepackage[utf8]{inputenc}
  \usepackage{textcomp} % provide euro and other symbols
\else % if luatex or xetex
  \usepackage{unicode-math} % this also loads fontspec
  \defaultfontfeatures{Scale=MatchLowercase}
  \defaultfontfeatures[\rmfamily]{Ligatures=TeX,Scale=1}
\fi
\usepackage{lmodern}
\ifPDFTeX\else
  % xetex/luatex font selection
\fi
% Use upquote if available, for straight quotes in verbatim environments
\IfFileExists{upquote.sty}{\usepackage{upquote}}{}
\IfFileExists{microtype.sty}{% use microtype if available
  \usepackage[]{microtype}
  \UseMicrotypeSet[protrusion]{basicmath} % disable protrusion for tt fonts
}{}
\makeatletter
\@ifundefined{KOMAClassName}{% if non-KOMA class
  \IfFileExists{parskip.sty}{%
    \usepackage{parskip}
  }{% else
    \setlength{\parindent}{0pt}
    \setlength{\parskip}{6pt plus 2pt minus 1pt}}
}{% if KOMA class
  \KOMAoptions{parskip=half}}
\makeatother
\setlength{\emergencystretch}{3em} % prevent overfull lines
\providecommand{\tightlist}{%
  \setlength{\itemsep}{0pt}\setlength{\parskip}{0pt}}
\usepackage{bookmark}
\IfFileExists{xurl.sty}{\usepackage{xurl}}{} % add URL line breaks if available
\urlstyle{same}
\hypersetup{
  hidelinks,
  pdfcreator={LaTeX via pandoc}}

\author{}
\date{}

\begin{document}

PONTIFÍCIA UNIVERSIDADE CATÓLICA DO RIO GRANDE DO SUL

ESCOLA POLITÉCNICA

Curso de BACHARELADO EM ENGENHARIA DE SOFTWARE

AGES - agência experimental de engenharia de software

HENRIQUE PINHO VALENTE SCHULTZ

memorial de atuação na agência experimental de engenharia de software --
período 09/08/2024 a 21/09/2026

AGES III

Porto Alegre

2026

\textbf{\hfill\break
}

\section{RESUMO}\label{resumo}

Este relatório apresenta minha trajetória nas experiências AGES I, AGES
II e AGES III (em andamento) evidenciando minha evolução técnica,
acadêmica e profissional ao longo do desenvolvimento de três projetos
distintos.

No AGES I, participei do projeto VIR -- Visual Interactive Reading, onde
vivenciei o processo inicial de um sistema: levantamento de requisitos,
criação de protótipos, organização de ambiente, desenvolvimento das
primeiras funcionalidades em React e integração entre frontend, backend
e banco de dados.

No AGES II, atuei no projeto Pró-Mata, aprofundando meu domínio em
desenvolvimento web moderno com React, TypeScript, Tailwind e TanStack
Router, além da criação de componentes reutilizáveis, fluxos de
navegação, design system inicial e estudos de validação e testes.

Já no AGES III, estou atuando no projeto FinUp com foco em arquitetura,
DevOps e infraestrutura. Tenho contribuido com a definição da
arquitetura do sistema, a estruturação dos repositórios no GitHub, o
estudo de serviços AWS, a configuração de pipelines de CI/CD, a adoção
da ferramenta Railforge para padronização do desenvolvimento e a
realização de code reviews. Essa experiência tem ampliado minhas
habilidades em decisões arquiteturais, containerização, integração
contínua, organização e comunicação de equipes, organização de código, e
entendimento de produto, contribuindo muito para meu crescimento como
desenvolvedor.

\textbf{Palavras-chave:} AGES; Engenharia de Software; Prototipação;
Desenvolvimento Web.

\textbf{\hfill\break
}

\section{LIsta de ilustrações}\label{lista-de-ilustrauxe7uxf5es}

Figura 1 - Foto da Equipe - VIR

Figura 2 - Foto da Equipe - PRO-MATA

\textbf{Figura 3 - Foto da Equipe - FinUp}

\section{\texorpdfstring{\hfill\break
}{ }}\label{section}

\section{SUMÁRIO}\label{sumuxe1rio}

\hyperref[apresentauxe7uxe3o-da-trajetuxf3ria-do-aluno]{1 - APRESENTAÇÃO
DA TRAJETÓRIA DO ALUNO} 5

\hyperref[_Toc45627142]{2 - PROJETOS AGES I - ``VIR -- Visual
Interactive Reading''} 7

\hyperref[introduuxe7uxe3o]{2.1 Introdução} 7

\hyperref[_Toc45627144]{2.2 Desenvolvimento do Projeto} 7

\hyperref[atividades-desempenhadas-pelo-aluno-no-projeto]{2.3 Atividades
desempenhadas pelo aluno no projeto} 9

\hyperref[conclusuxe3o]{2.4. Conclusão} 12

\hyperref[_Toc45627142]{3 - PROJETOS AGES II - ``PRO-MATA ''} 15

\hyperref[introduuxe7uxe3o]{3.1 Introdução} 15

\hyperref[_Toc45627144]{3.2 Desenvolvimento do Projeto} 15

\hyperref[atividades-desempenhadas-pelo-aluno-no-projeto]{3.3 Atividades
desempenhadas pelo aluno no projeto} 18

\hyperref[conclusuxe3o]{3.4. Conclusão} 23

\hyperref[_Toc45627142]{4 - PROJETOS AGES III - ``FINUP ''} 25

\hyperref[introduuxe7uxe3o]{4.1 Introdução} 25

\hyperref[_Toc45627144]{4.2 Desenvolvimento do Projeto} 25

\hyperref[atividades-desempenhadas-pelo-aluno-no-projeto]{4.3 Atividades
desempenhadas pelo aluno no projeto} 28

\hyperref[referuxeancias]{REFERÊNCIAS} 30

\section{1 - APRESENTAÇÃO DA TRAJETÓRIA DO
ALUNO}\label{apresentauxe7uxe3o-da-trajetuxf3ria-do-aluno}

Desde cedo, minha curiosidade esteve ligada à criação. Antes mesmo de
pensar em Engenharia de Software, eu explorava ferramentas como
Photoshop, Premiere, Blender e Animate para entender como imagens,
vídeos e animações ganhavam vida. Essa familiaridade com construção
visual despertou uma pergunta que se tornaria decisiva na minha
trajetória: \emph{``como o software por trás disso tudo funciona?''}.

Durante a quarentena, essa inquietação se tornou um ponto de virada.
Tranquei Economia, mergulhei por um período na criação de conteúdo e,
refletindo sobre futuro e propósito, percebi que minha motivação estava
menos em comunicar ideias e mais em construir soluções\textbf{.} Foi a
partir desse momento que decidi migrar definitivamente para TI e iniciar
meu caminho na Engenharia de Software.

Desde o início da graduação, busquei equilíbrio entre técnica,
organização e criatividade. No meu dia a dia, adotei ferramentas como
Notion, Miro, Trello, Forest e automações com Alexa, o que estruturou o
meu jeito de aprender: iterativo, visual e orientado a processos. Ao
longo do curso, fui desenvolvendo a convicção de que queria me tornar
full stack, capaz de transformar ideias em produtos completos.

Minha trajetória profissional começou a ganhar forma no estágio de dois
anos no Palácio da Justiça, experiência que, embora baseada em
ferramentas no-code para construção de aplicações, foi extremamente rica
em termos de maturidade profissional. Mesmo sem lidar diretamente com
código tradicional, desenvolvi habilidades fundamentais para a área de
TI: organização, comunicação, entendimento de requisitos, trabalho
colaborativo e responsabilidade com entregas. Além disso, tive a
oportunidade de participar de processos internos relevantes, incluindo a
condução de entrevistas para novos candidatos à vaga de TI, o que
ampliou minha visão sobre cultura de equipe e critérios de seleção. Essa
vivência consolidou minhas soft skills e me deu segurança para atuar em
ambientes reais de desenvolvimento. Ao final desse ciclo, fui aprovado
na Dell IT Academy, que expandiu minha visão sobre arquitetura,
qualidade e práticas modernas de desenvolvimento --- conhecimentos que
se tornaram essenciais para minha evolução dentro da AGES.

A transição do AGES I para o AGES II marcou um ponto significativo de
evolução. Cheguei no novo ciclo mais maduro tecnicamente e com uma visão
mais clara de produto: antes de codar, eu precisava entender\textbf{,}
esboçar\textbf{,} prototipar. Esse pensamento guiou meu processo no
projeto Pró-Mata, onde passei a validar fluxos no Figma, aprimorar
interações e antecipar decisões de UX antes de escrever qualquer linha
de código. A prática de ``prototipar primeiro, codificar depois''
reduziu retrabalho, alinhou expectativas com cliente e tornou meu
desenvolvimento mais intencional.

Do ponto de vista técnico, foi no AGES II que descobri um conjunto de
tecnologias modernas que facilitaram meu desenvolvimento e ampliaram
minha compreensão sobre como um frontend profissional é estruturado.
Passei a utilizar ferramentas como React aliado ao TypeScript, Tailwind
CSS, TanStack Router, Zustand, React Hook Form e Zod, entendendo como
cada uma contribui para tornar o código mais organizado e previsível.
Nesse processo, também amadureci meu uso de Git/GitHub e consolidei
hábitos de organização e comunicação que tornaram meu trabalho mais
estruturado e profissional.

Ao mesmo tempo, meu crescimento não foi apenas técnico. A participação
ativa nas discussões, a segurança para propor ideias, a abertura para
feedback e a melhora na comunicação com o time marcaram meu avanço em
soft skills. Passei a enxergar o desenvolvimento não só como código, mas
como colaboração, clareza e tomada de decisão\textbf{.} Meu papel deixou
de ser apenas entregar tarefas e passou a envolver visão de produto,
entendimento de fluxo, coerência visual e contribuição para um design
system inicial.

Hoje, encerro minha AGES II com uma base muito mais sólida: rotas,
componentes, padrões visuais e arquitetura de estado definidos;
compreensão mais profunda da relação entre design e código; e uma visão
mais completa de como um sistema nasce no papel e ganha forma na
prática.

Essa evolução acadêmica, técnica e pessoal define minha jornada até
aqui. A AGES tem sido não apenas um espaço de prática, mas um marco
contínuo na construção da minha identidade como desenvolvedor: mais
consciente, mais consistente e mais preparado para transformar ideias em
soluções reais.

\protect\phantomsection\label{_Toc45627142}{}

\section{2 - PROJETOS AGES I - ``VIR -- Visual Interactive
Reading''}\label{projetos-ages-i---vir-visual-interactive-reading}

\subsection{2.1 Introdução}\label{introduuxe7uxe3o}

\protect\phantomsection\label{_Toc45627144}{}O projeto "VIR -- Visual
Interactive Reading" busca desenvolver um aplicativo interativo de
leitura para jovens de 12 a 15 anos, alinhado às competências visuais da
geração Alpha. O objetivo é oferecer uma experiência que combina textos
de autores clássicos e contemporâneos com animações, estimulando a
criatividade visual e engajando leitores em formação.

O VIR surge da necessidade identificada em escolas do ensino
fundamental, que enfrentam desafios na formação de leitores engajados.
Baseado em pesquisas acadêmicas e entrevistas com educadores, o
aplicativo contará com funcionalidades como login personalizado, estante
digital, controle de leitura individualizada e recomendação de obras por
faixa etária. Essa abordagem busca complementar a leitura tradicional,
tornando-a mais atrativa e interativa para o público jovem.

\subsection{2.2 Desenvolvimento do
Projeto}\label{desenvolvimento-do-projeto}

\textbf{Stakeholders do Projeto:}

\begin{itemize}
\item
  Hermílio Santos (Professor da Escola de Humanidades - Sociologia -
  PUCRS)
\item
  Paula Vianna (Mestranda em Sociologia - Escola de Humanidades - PUCRS)
\end{itemize}

\textbf{Período de Execução}: Segundo semestre de 2024.

Figura - Foto do time

\begin{enumerate}
\def\labelenumi{\arabic{enumi}.}
\item
  \textbf{Repositório do código Fonte do Projeto}
\end{enumerate}

\url{https://tools.ages.pucrs.br/VIR/vir-wiki/-/tree/main}

O código fonte está hospedado no GitLab, onde todo o desenvolvimento
colaborativo é gerenciado.

\begin{enumerate}
\def\labelenumi{\arabic{enumi}.}
\setcounter{enumi}{1}
\item
  \textbf{Banco de Dados utilizado}
\end{enumerate}

\url{https://tools.ages.pucrs.br/VIR/vir-wiki/-/wikis/banco_dados}

\begin{quote}
Diagrama das tabelas do banco de dados, incluindo usuários, livros,
capítulos, páginas, vídeos e progresso de leitura, ilustrando suas
relações e funcionalidades no projeto VIR.
\end{quote}

\begin{enumerate}
\def\labelenumi{\arabic{enumi}.}
\setcounter{enumi}{2}
\item
  \textbf{Arquitetura utilizada}
\end{enumerate}

\url{https://tools.ages.pucrs.br/VIR/vir-wiki/-/wikis/arquitetura}

Arquitetura utliza React no frontend, Flask na API backend e PostgreSQL
como banco de dados.

\begin{enumerate}
\def\labelenumi{\arabic{enumi}.}
\setcounter{enumi}{3}
\item
  \textbf{Protótipos das telas desenvolvidas}
\end{enumerate}

\url{https://tools.ages.pucrs.br/VIR/vir-wiki/-/wikis/design_mockups}

Os mockups do projeto foram feitos no Figma, exibindo o design e fluxo
das interfaces.

\begin{enumerate}
\def\labelenumi{\arabic{enumi}.}
\setcounter{enumi}{4}
\item
  \textbf{Tecnologias Utilizadas}
\end{enumerate}

\url{https://tools.ages.pucrs.br/VIR/vir-wiki/-/wikis/configuracao}

React (Frontend): Uma biblioteca JavaScript para construir interfaces de
usuário dinâmicas e modernas, permitindo a criação de componentes
reutilizáveis e otimizando a atualização de interfaces por meio do
virtual DOM.

Flask (Backend): Microframework web para Python que permite a criação de
APIs RESTful e aplicações web de forma simples e flexível.

Python: Linguagem de programação usada no backend devido à sua sintaxe
clara e suporte extensivo a bibliotecas.

PostgreSQL: Sistema de gerenciamento de banco de dados relacional,
conhecido por sua robustez e conformidade com padrões SQL.

\subsection{2.3 Atividades desempenhadas pelo aluno no
projeto}\label{atividades-desempenhadas-pelo-aluno-no-projeto}

\textbf{2.3.1 Sprint 0}

Durante a Sprint 0, meu foco foi compreender os requisitos do cliente e
dar os primeiros passos no desenvolvimento do projeto. As atividades
foram estruturadas para garantir que as bases do projeto estivessem bem
estabelecidas.

\textbf{Atividades Desempenhadas:}

\begin{itemize}
\item
  Preparação para a Reunião Inicial: Desenvolvi perguntas para a reunião
  com o cliente, focando em esclarecer os requisitos, funcionalidades,
  design e aspectos técnicos do sistema.
\item
  Desenvolvimento de Mockups no Figma: Trabalhei na criação de mockups,
  com ênfase no footer do aplicativo. Testei diferentes versões para
  aprimorar a interface e explorei ideias para o scroll da
  funcionalidade de leitura, inspirando-me no estilo de scroll do
  TikTok.
\item
  Configuração do Ambiente de Desenvolvimento: Clonei o repositório do
  GitLab para minha máquina local, ajustei o ambiente, instalei Node.js
  e suas dependências. Fiz a configuração inicial necessária para
  garantir o funcionamento do projeto.
\item
  Revisão de Git e Shell: Para melhorar o fluxo de desenvolvimento,
  revisei o uso do Git e comandos do shell, fundamentais para a
  organização e colaboração no projeto.
\item
  Correção de Vulnerabilidades: Durante a instalação do Node.js,
  encontrei vulnerabilidades que foram corrigidas. Com a ajuda dos
  alunos do AGES IV, entendi que essas vulnerabilidades não impactariam
  o desenvolvimento inicial.
\end{itemize}

\textbf{Ferramentas Utilizadas:}

\begin{itemize}
\item
  Figma: Utilizado para a criação dos mockups e testes de interface.
\item
  GitLab: Ferramenta de versionamento usada para clonar e gerenciar o
  código-fonte do projeto.
\item
  Node.js: Instalado como base para o backend do projeto, com as
  dependências necessárias configuradas.
\item
  Comandos de Git e Shell: Revisados e utilizados para a gestão do
  repositório e ambiente local.
\end{itemize}

\textbf{Contribuições no Desenvolvimento:\\
}\strut \\
Minhas principais contribuições nesta sprint foram a criação dos mockups
iniciais no Figma, com foco em uma interface intuitiva, e a configuração
do ambiente de desenvolvimento. A preparação detalhada para a reunião
inicial com o cliente também ajudou a garantir que as expectativas e
requisitos do sistema fossem claramente definidos.

\textbf{2.3.2 Sprint 1}

Na Sprint 1, o foco foi o desenvolvimento das funcionalidades de
navegação e layout do aplicativo, com destaque para o componente de
footer e o fluxo de leitura, além do estudo contínuo de tecnologias
essenciais para o projeto.

\textbf{Atividades Desempenhadas:}

\begin{itemize}
\item
  Estudo de React: Aprofundei meus conhecimentos em React para
  desenvolver as funcionalidades de navegação do aplicativo. Isso
  incluiu a criação de componentes e o gerenciamento de estado,
  essenciais para o funcionamento do sistema.
\item
  Criação do Footer Component: Desenvolvi o componente de footer na tela
  de estante/biblioteca, garantindo uma navegação eficiente. Fiz ajustes
  no layout utilizando CSS, centralizando e estilizando ícones SVG para
  uma apresentação mais organizada e responsiva.
\item
  Atualização e Configuração do Ambiente: Reconfigurei o ambiente de
  desenvolvimento e criei uma nova branch específica para o
  desenvolvimento do Footer Component. Isso ajudou a manter o controle
  das modificações e facilitou o trabalho colaborativo.
\item
  Implementação Inicial das Rotas: Iniciei a implementação das rotas de
  navegação entre as páginas do aplicativo através dos botões do footer,
  mas não consegui finalizá-las devido à carga de provas e trabalhos
  acadêmicos acumulados.
\item
  Acompanhamento do Time: Embora meu progresso tenha sido limitado,
  acompanhei o desenvolvimento das tarefas da minha squad, garantindo
  que o andamento do projeto continuasse conforme o planejado.
\end{itemize}

\textbf{Ferramentas Utilizadas:}

\begin{itemize}
\item
  React: Utilizado para criar o componente de footer e gerenciar a
  navegação entre as páginas do aplicativo.
\item
  CSS e SVG: CSS foi aplicado para estilizar o layout e centralizar os
  ícones SVG, melhorando o controle visual e responsividade.
\item
  GitLab: Utilizado para gerenciar o repositório do projeto e criar a
  branch dedicada ao Footer Component.
\item
  Node.js: Mantido como base do ambiente de desenvolvimento para
  garantir o funcionamento adequado da aplicação.
\end{itemize}

\textbf{Contribuições no Desenvolvimento:\\
}\strut \\
Minha principal contribuição nesta sprint foi o desenvolvimento do
Footer Component e o início da implementação das rotas de navegação,
essenciais para o fluxo do aplicativo. Apesar das limitações impostas
pela carga acadêmica, garanti que o layout e a estrutura do footer
estivessem finalizados e funcionais, além de contribuir para o andamento
da squad ao acompanhar as tarefas do time.

\textbf{2.3.3 Sprint 2}

Na Sprint 2, o foco foi o desenvolvimento da página de login do VIR,
assegurando funcionalidade, integração com o backend e design intuitivo.
As tarefas incluíram desde a criação e estilização do componente até a
resolução de problemas encontrados durante o desenvolvimento.

\textbf{Atividades Desempenhadas:}

\begin{itemize}
\item
  Criação de uma nova branch: Iniciei o desenvolvimento criando a branch
  feature/Student/Login/Create-Login-Page a partir da branch dev,
  garantindo organização e controle das modificações.
\item
  Desenvolvimento do Componente Form: Desenvolvi e estilizei o
  componente Form, que incluiu os campos de input para "Email" e
  "Senha", com funcionalidade de censura no campo de senha.
\item
  Criação da Página de Login: Criei e estilizei a página de login,
  integrando o símbolo do VIR e o componente Form, garantindo uma
  interface funcional e intuitiva.
\item
  Configuração e Integração: Implementei o envio de um arquivo JSON
  contendo email e senha ao endpoint correspondente, assegurando
  funcionalidade básica de autenticação.
\item
  Controle de Versionamento: Realizei commit das alterações, rebase com
  a branch dev para manter as atualizações e push para o repositório
  remoto.
\item
  Resolução de Problemas no GitLab: Ao fazer um pull request para a
  branch errada (main), corrigi rapidamente com a ajuda dos colegas do
  AGES IV, garantindo a manutenção do fluxo de trabalho
\end{itemize}

\textbf{Ferramentas Utilizadas:}

\begin{itemize}
\item
  GitLab: Utilizado para criação da branch, versionamento do código e
  push das alterações.
\item
  Node.js: Mantido no ambiente de desenvolvimento para suporte ao envio
  do JSON.
\item
  CSS: Usado para estilizar a interface e aprimorar o layout da página
  de login.
\item
  React: Empregado no desenvolvimento do componente Form e na integração
  da página de login.
\end{itemize}

\textbf{Contribuições no Desenvolvimento:}

Minha principal contribuição nesta sprint foi a criação e estilização da
página de login, incluindo o componente Form e a integração com o
backend para envio do JSON. Além disso, a resolução de problemas
relacionados ao GitLab contribuiu para o fluxo de trabalho da equipe,
garantindo organização e eficiência.

\textbf{2.3.4 Sprint 3}

Na Sprint 3, o foco principal foi o desenvolvimento do SchoolFormModal,
com atenção especial à sua estrutura, estilização e funcionalidade, além
de estudos relacionados ao gerenciamento de estado no React.

\textbf{Atividades Desempenhadas:}

\begin{itemize}
\item
  Estruturação do SchoolFormModal: Desenvolvi o modal com header, footer
  e ícone "X" (importado como PNG), além dos botões "Cancelar" e
  "Confirmar".
\item
  Organização de Arquivos CSS: Criei arquivos CSS individuais para cada
  componente do SchoolFormModal, permitindo maior modularidade e
  manutenção.
\item
  Estilização de Botões: Estilizei o botão "Confirmar" com ajustes
  visuais, como cor, sombra e comportamentos específicos para os estados
  hover e click.
\item
  Criação de Componentes Reutilizáveis: Organizei os botões "Cancelar" e
  "Confirmar" como componentes individuais na pasta adminCommons.
\item
  Integração Visual: Adicionei o SchoolFormModal à página da escola
  (escolaPage) para acompanhar o desenvolvimento visual no contexto do
  sistema.
\item
  Estudo de React: Realizei estudos aprofundados sobre o hook useState,
  compreendendo melhor o gerenciamento de estado em componentes
  funcionais.
\item
  Estilização de Inputs: Estilizei os campos de input do modal e revisei
  a estilização geral, adicionando estados hover, focus e pressed aos
  elementos.
\end{itemize}

\textbf{Ferramentas Utilizadas:}

\begin{itemize}
\item
  React: Utilizado para desenvolver o modal, criar componentes
  individuais e implementar o gerenciamento de estado com o hook
  useState.
\item
  CSS: Aplicado na estilização de botões e inputs, com foco em
  modularidade e estados dinâmicos.
\item
  GitLab: Mantido como ferramenta de versionamento para gerenciar as
  alterações no projeto.
\end{itemize}

\textbf{Contribuições no Desenvolvimento:}

Minha principal contribuição nesta sprint foi o desenvolvimento completo
do SchoolFormModal\textbf{,} abrangendo estrutura, estilização e
integração visual na página da escola. Além disso, os estudos sobre o
hook useState aumentaram minha eficiência no gerenciamento de estados em
componentes React.

\textbf{2.3.5 Sprint 4}

Na Sprint 4, o foco foi o desenvolvimento e finalização das
funcionalidades relacionadas à gestão de alunos, incluindo os
componentes Student Modal Form, Student List Component e Student Page,
além de melhorias no controle de estado e correção de problemas
técnicos.

\textbf{Atividades Desempenhadas:}

\begin{itemize}
\item
  Correção do Controle de Estado: Resolvi problemas no controle de
  abertura (isOpen) do modal no componente AlunoPage, ajustando o uso de
  props entre componentes pai e filho.
\item
  Estudo sobre CSS: Realizei estudos sobre propriedades de display no
  CSS, como Grid, Flexbox e alinhamentos (align-items, align-content,
  justify-items, justify-content), aprofundando o entendimento sobre as
  melhores práticas de estilização.
\item
  Desenvolvimento de Componentes: Concluí o desenvolvimento e a
  estilização dos componentes Student Page, Student Modal Form e Student
  List Component, garantindo funcionalidade e integração visual.
\item
  Push e Merge Request: Realizei push e merge request da página de
  aluno, garantindo a sincronização das alterações no repositório
  remoto.
\item
  Investigação de Falhas: Identifiquei e tentei corrigir falhas na
  pipeline devido a branch desatualizada, realizando rebase para
  resolver os conflitos.
\item
  Implementação de Endpoints: Iniciei a implementação dos endpoints da
  página de aluno, enfrentando obstáculos técnicos que ainda precisam de
  resolução.
\end{itemize}

\textbf{Ferramentas Utilizadas:}

\begin{itemize}
\item
  React: Utilizado para corrigir o controle de estado e desenvolver os
  componentes relacionados à página de aluno.
\item
  CSS: Aplicado para estilização dos componentes, com foco no uso
  prático de Grid e Flexbox.
\item
  GitLab: Usado para realizar o push, rebase e merge request das
  alterações no repositório.
\item
  Node.js: Base para a tentativa de implementação dos endpoints.
\end{itemize}

\textbf{Contribuições no Desenvolvimento:}

Minha principal contribuição nesta sprint foi o desenvolvimento e
finalização da Student Page\textbf{,} Student Modal Form e Student List
Component, essenciais para a gestão de alunos. Além disso, resolvi
problemas técnicos relacionados ao controle de estado no modal e
contribuí para a estabilização do fluxo de trabalho no repositório.

\subsection{2.4. Conclusão}\label{conclusuxe3o}

Este projeto se destacou como uma experiência transformadora na minha
trajetória acadêmica, pois me proporcionou uma visão ampla e prática do
fluxo de desenvolvimento de um projeto real. Desde a definição de
requisitos até a entrega de funcionalidades, pude compreender como as
diferentes etapas de um projeto se interconectam, exigindo planejamento,
colaboração e atenção aos detalhes.

O desenvolvimento do VIR exigiu que eu aprofundasse meu entendimento
sobre a estruturação de um projeto como um todo, incluindo a integração
entre o front-end, o back-end e o banco de dados. Trabalhar em um
ambiente dinâmico me permitiu visualizar como essas áreas se comunicam,
como gerenciar estados e como organizar componentes modulares e
reutilizáveis de forma eficiente. Além disso, o uso de ferramentas como
React, GitLab e CSS reforçou a importância de boas práticas de
desenvolvimento e versionamento.

Essa experiência também destacou a relevância de metodologias ágeis para
organizar o trabalho em equipe e garantir entregas contínuas. Aprendi a
lidar com problemas técnicos, como falhas na pipeline e dificuldades de
integração, e a importância de alinhar expectativas com a equipe e
stakeholders. Esse aprendizado me preparou para lidar melhor com a
complexidade de projetos futuros, ampliando minha confiança e
competência técnica.

O desenvolvimento do projeto também me fez perceber a importância de
alinhar a teoria com a prática, exigindo esforço para aplicar
conhecimentos de programação, design e gerenciamento em um contexto
real. No futuro, pretendo aproveitar esse aprendizado para me tornar um
profissional mais completo, capaz de colaborar em projetos que demandem
integração entre diversas áreas e habilidades.

\section{3 - PROJETOS AGES II -
``PRO-MATA''}\label{projetos-ages-ii---pro-mata}

\subsection{3.1 Introdução}\label{introduuxe7uxe3o-1}

O Centro de Pesquisas e Proteção da Natureza (CPPN) Pró-Mata desempenha
um papel essencial na preservação ambiental e também na promoção de
atividades acadêmicas, turísticas e culturais. Para atender melhor esse
público diverso, surgiu a necessidade de criar uma plataforma digital
que centralize e facilite o acesso aos serviços oferecidos pelo centro.

A iniciativa busca modernizar e tornar mais acessíveis processos que
antes eram manuais ou fragmentados, como reservas de hospedagem,
agendamento de atividades e organização de eventos. Com a plataforma,
espera-se proporcionar mais praticidade e transparência aos visitantes,
ao mesmo tempo em que se fortalece a gestão interna e se melhora a
experiência de todos que têm contato com o Pró-Mata --- sejam
estudantes, pesquisadores, turistas ou instituições parceiras.

Mais do que uma solução tecnológica, este projeto representa um passo
importante para aproximar o centro de sua comunidade, ampliando seu
alcance e reforçando sua missão de unir conservação, educação e turismo
sustentável.

\subsection{3.2 Desenvolvimento do
Projeto.}\label{desenvolvimento-do-projeto.}

\textbf{Stakeholder do Projeto:}

\begin{itemize}
\item
  Prof. Augusto Mussi Alvim
\end{itemize}

\textbf{Período de Execução}: Segundo semestre de 2025.

Figura - Foto do time

\begin{enumerate}
\def\labelenumi{\arabic{enumi}.}
\item
  \textbf{Repositório do código Fonte do Projeto}
\end{enumerate}

\url{https://tools.ages.pucrs.br/pro-mata}

O código fonte está hospedado no GitLab, onde todo o desenvolvimento
colaborativo é gerenciado.

\begin{enumerate}
\def\labelenumi{\arabic{enumi}.}
\setcounter{enumi}{1}
\item
  \textbf{Banco de Dados utilizado}
\end{enumerate}

\url{https://tools.ages.pucrs.br/pro-mata/wiki/-/wikis/Banco\%20de\%20Dados}

\begin{quote}
Diagrama das tabelas do banco de dados.

\textbf{Desenvolvimento}: Um container PostgreSQL simples que roda
localmente. O Prisma ORM gerencia as migrações automaticamente.

\textbf{Produção}: PostgreSQL com replicação master-slave, backups
automáticos criptografados, e monitoramento futuro com Grafana e
PostgreSQL exporter.
\end{quote}

\begin{enumerate}
\def\labelenumi{\arabic{enumi}.}
\setcounter{enumi}{2}
\item
  \textbf{Arquitetura utilizada}
\end{enumerate}

https://tools.ages.pucrs.br/pro-mata/wiki/-/wikis/Arquitetura

Baseando-se no que foi planejado para o nosso banco de dados,
levantamento dos requisitos e o entendimento geral do time sobre o
projeto a ser desenvolvido, optamos por uma arquitetura de
Cliente-Servidor, onde há uma separação entre o frontend e backend.

\begin{enumerate}
\def\labelenumi{\arabic{enumi}.}
\setcounter{enumi}{3}
\item
  \textbf{Protótipos das telas desenvolvidas}
\end{enumerate}

https://tools.ages.pucrs.br/pro-mata/wiki/-/wikis/Design

Os mockups do projeto foram feitos no Figma, exibindo o design e fluxo
das interfaces.

\begin{enumerate}
\def\labelenumi{\arabic{enumi}.}
\setcounter{enumi}{4}
\item
  \textbf{Tecnologias Utilizadas}
\end{enumerate}

\begin{quote}
https://tools.ages.pucrs.br/pro-mata/wiki/-/wikis/Arquitetura

\textbf{React 19}: Construção de interfaces declarativas com componentes
reutilizáveis.

\textbf{TypeScript:} Tipagem estática (configuração em tsconfig.json,
moduleResolution: "bundler", baseUrl, paths @/*).

\textbf{Vite 6} --- Dev server e build otimizados; plugin
\textbf{@tanstack/router-plugin} com autoCodeSplitting e geração de
src/routeTree.gen.ts; plugin \textbf{@tailwindcss/vite}.

\textbf{Tailwind CSS (v4)} --- Utilitários CSS, tokens/temas em
src/styles/globals.css (variações claro/escuro).

\textbf{TanStack Router} --- Roteamento file-based em src/routes/**, com
beforeLoad/guards, code-splitting e devtools.

\textbf{TanStack Query} --- Cache e invalidação de dados; QueryClient
injetado no contexto do Router.

\textbf{shadcn/ui + Radix UI} --- Componentes acessíveis (dialog,
popover, select, table, sheet, etc.) com variantes (CVA).

\textbf{Zod + React Hook Form} --- Esquemas de validação tipados e
formulários controlados.

\textbf{Zustand} --- Estado leve e local (stores de UI e filtros).

\textbf{Axios} --- Cliente HTTP com \textbf{interceptor} de autenticação
(Bearer).

\textbf{ESLint/Prettier} --- Lint e formatação padronizada.

\textbf{Vitest} --- Testes unitários com run, watch e UI.

\textbf{Node.js / npm} --- Toolchain e scripts de desenvolvimento.
\end{quote}

\subsection{3.3 Atividades desempenhadas pelo aluno no
projeto}\label{atividades-desempenhadas-pelo-aluno-no-projeto-1}

\textbf{3.3.1 Sprint 0}

A Sprint 0 foi dedicada a alinhar a visão do produto com protótipos
navegáveis no Figma e a colocar no ar os primeiros componentes do
frontend, organizando o fluxo de telas e preparando a base visual para
as próximas sprints.

\textbf{Atividades Desempenhadas:}

\begin{itemize}
\item
  Desenvolvimento no Figma: ajustes de mockups, criação de cards
  compactos, animações no Prototype e botão de idioma.
\item
  Definição do \textbf{System Flow} (user flow/screen flow) para
  orientar navegação e decisões de UI.
\item
  Implementação do \textbf{CardStatus} (enum de estados, ícones SVG,
  fallback para status desconhecido e estilos com Tailwind).
\item
  Desenvolvimento do \textbf{CardsInfoOnHover} (branch \#44): estrutura
  HTML, estilos com Tailwind, lógica de hover com useState, integração
  de imagens mock, posicionamento na Home; revisão de tipagem, ajustes
  com TanStack Router e \textbf{PR} entregue.
\item
  Ambientação/estudo do stack do frontend (React, TypeScript, Tailwind,
  TanStack Router, Zustand, React Hook Form + Zod, JWT) e ajustes de
  ambiente (VS Code, PostCSS, ícones).
\end{itemize}

\textbf{Ferramentas Utilizadas:}

\begin{itemize}
\item
  \textbf{Figma} (mockups e Prototype)
\item
  \textbf{React} + \textbf{TypeScript}
\item
  \textbf{Tailwind CSS}
\item
  \textbf{TanStack Router}
\item
  \textbf{Zustand}
\item
  \textbf{React Hook Form} + \textbf{Zod}
\item
  \textbf{Git/GitHub}, \textbf{VS Code}, PostCSS, SVGs
\end{itemize}

\textbf{Contribuições no desenvolvimento:}

\begin{itemize}
\item
  Protótipos navegáveis que alinham expectativas e guiam o backlog.
\item
  Componentes prontos e reutilizáveis (\textbf{CardStatus} e
  \textbf{CardsInfoOnHover}), já padronizados visualmente.
\item
  Base de estilos, ícones e estados definida para acelerar novas telas.
\item
  PRs documentado que facilita revisão e evolução do código.
\end{itemize}

\textbf{3.3.2 Sprint 1}

A Sprint 1 foi voltada a transformar os protótipos em telas navegáveis e
consolidar a base do frontend. Houve foco em roteamento, componentes
reutilizáveis (Carrossel, Cards) e padronização visual. Apesar da curva
de aprendizado com TanStack Router, tipagem em TS e detalhes de
Tailwind, o time encerrou a sprint com a Home estruturada e
componentes-chave prontos para evoluir.

\textbf{Atividades Desempenhadas:}

\begin{itemize}
\item
  Estudo dirigido da stack (React + TypeScript, Tailwind, TanStack
  Router, React Hook Form + Zod, JWT) e ajustes de ambiente (VS Code,
  PostCSS, Montserrat, organização de assets/mocks).
\item
  \textbf{CardsInfoOnHover (\#44):} refinado e entregue via PR (tipagem
  revisada, hover com estado, integração de imagens mock e
  posicionamento na Home).
\item
  \textbf{CarouselComponent (\#43):} base implementada com navegação por
  setas, miniaturas e exibição de 2 imagens por vez; controle de índices
  para evitar overflow.
\item
  \textbf{Home (\#42):} estrutura inicial com Header/Footer, CTA e
  imagem de destaque, alinhada aos protótipos.
\item
  \textbf{ExperienceCard (\#78):} esqueleto criado com props iniciais e
  estados de loading usando react-spinners.
\end{itemize}

\textbf{Ferramentas Utilizadas:}

\begin{itemize}
\item
  \textbf{React} + \textbf{TypeScript}, \textbf{Tailwind CSS},
  \textbf{TanStack Router}, \textbf{React Hook Form} + \textbf{Zod},
  \textbf{JWT} (estudo/integração inicial).
\item
  \textbf{Figma} (referência de layout e fluxo).
\item
  \textbf{react-spinners}, \textbf{SVGs reutilizáveis} (ex.: ArrowSVG).
\item
  \textbf{VS Code}, \textbf{Git/GitHub}, \textbf{PostCSS}; fonte
  \textbf{Montserrat} configurada; \textbf{public/} para mocks e
  imagens.
\end{itemize}

\textbf{Contribuições no desenvolvimento:}

\begin{itemize}
\item
  Home navegável e padronizada como base para novas seções.
\item
  Carrossel reutilizável (\#43) com UX funcional e regras de borda
  implementadas.
\item
  CardsInfoOnHover (\#44) maduro e integrado, com PR documentado.
\item
  Base do ExperienceCard preparada para conectar a dados reais (props,
  loading).
\item
  Padronização de ícones, diretórios (components/ui/svgs/) e uso
  consistente de cores do global.css, aumentando a coesão visual do
  projeto.
\end{itemize}

\textbf{3.3.3 Sprint 2}

A Sprint 2 teve foco no desenvolvimento do componente ExperienceCard,
refinando o design system da tela de reservas, ajustando detalhes
visuais e preparando o componente para suportar diferentes tipos de
experiências, já alinhado à arquitetura fullstack pensada para o
projeto.

\textbf{Atividades Desempenhadas:}

\begin{itemize}
\item
  Criação da branch feature/experience-card a partir da develop.\\
  Instalação e configuração do react-spinners para estados de loading no
  card.
\item
  Implementação do componente ExperienceCard com layout dinâmico
  conforme experience.category (trilha, hospedagem, laboratório,
  evento).
\item
  Estruturação do layout interno em duas colunas, com imagem, labels,
  descrição e botão ``Adicionar ao carrinho''.
\item
  Criação de mocks para as quatro categorias e exibição em grid na rota
  /reserve.
\item
  Aplicação de cores padronizadas via global.css e criação da variável
  -\/-cor-card-labels.
\item
  Ajustes visuais (sombra inset, scroll invisível para descrições
  longas, padding/margin, z-index).
\item
  Configuração de labels com ícones dinâmicos via Lucide React.
\item
  Substituição de textos estáticos por tags de Locale, preparando o
  componente para internacionalização.
\item
  Revisão de PR com base em feedback técnico, remoção de
  imports/variáveis não utilizados.
\item
  Correção do CI ao limpar declarações não utilizadas em arquivos
  TypeScript.
\end{itemize}

\textbf{Ferramentas Utilizadas:}

\begin{itemize}
\item
  React + TypeScript, Tailwind CSS.
\item
  Lucide React (ícones), react-spinners.
\item
  TanStack Router (rota /reserve).
\item
  Git/GitHub (branching, PRs, code review).
\item
  global.css (tokens de cor e variáveis globais).
\end{itemize}

\textbf{Contribuições no desenvolvimento:}

\begin{itemize}
\item
  ExperienceCard funcional e flexível, pronto para receber dados reais e
  suportar diferentes categorias de experiência.
\item
  Refinamento do design system (cores globais, sombras, labels,
  comportamento de scroll), aumentando a consistência visual.
\item
  Preparação da tela de reservas com mocks e grid montado, facilitando a
  futura integração com backend.
\item
  Melhoria na qualidade do código e estabilidade do pipeline de CI, a
  partir da limpeza de imports e variáveis não utilizados.
\item
  Avanço na compreensão da arquitetura fullstack e do uso de TanStack
  Query para cache e invalidação de dados (estudo inicial).
\end{itemize}

\textbf{3.3.4 Sprint 3}

A Sprint 3 marcou o início efetivo da minha atuação em testes de
frontend, redirecionando parte do foco do desenvolvimento de componentes
para a criação de uma base de qualidade com testes unitários,
entendimento da infraestrutura de mocks e organização do ecossistema de
testes do projeto.

\textbf{Atividades Desempenhadas:}

\begin{itemize}
\item
  Participação em reunião com os AGES II para definição e divisão das
  tasks de testes.
\item
  Resolução de conflitos de merge entre branches do frontend
  relacionados à área de testes.
\item
  Estudo da estrutura de testes do projeto, compreendendo a separação
  entre src/test/msw (infraestrutura de mocks) e tests/unit (testes
  reais).
\item
  Aprofundamento no papel de cada ferramenta:

  \begin{itemize}
  \item
    Node.js como ambiente de execução.
  \item
    Vite como dev server e ferramenta de build.
  \item
    Vitest como test runner.
  \item
    React Testing Library para simulação de uso em jsdom.
  \end{itemize}
\item
  Criação do primeiro teste unitário real cobrindo o componente
  CardStatus (verificação de renderização de label conforme status).
\item
  Início da implementação dos testes do componente ExperienceCard,
  estruturando cenários iniciais de renderização.
\end{itemize}

\textbf{Ferramentas Utilizadas:}

\begin{itemize}
\item
  Vitest (test runner).
\item
  React Testing Library (renderização e queries).
\item
  MSW (infraestrutura de mocks, estudo da pasta src/test/msw).
\item
  Node.js, Vite (ambiente de execução dos testes).
\item
  Git/GitHub (resolução de conflitos e organização de branches).
\end{itemize}

\textbf{Contribuições no desenvolvimento:}

\begin{itemize}
\item
  Estrutura inicial de testes unitários aplicada em um componente já
  utilizado na Home (CardStatus).
\item
  Organização e compreensão da arquitetura de testes do projeto,
  facilitando a evolução futura de testes em outros componentes.
\item
  Suporte na resolução de conflitos de merge, contribuindo para a
  estabilidade da branch principal.
\item
  Primeiros passos na criação de testes para o ExperienceCard,
  preparando o terreno para aumentar a cobertura de testes dos
  componentes de reserva.
\end{itemize}

\textbf{3.3.5 Sprint 4}

A Sprint 4 deu continuidade ao trabalho iniciado em testes, consolidando
o aprendizado em Vitest e Testing Library, expandindo a cobertura de
componentes críticos (ExperienceCard e Tela de Reservas) e colaborando
com a estabilidade do código na reta final da disciplina.

\textbf{Atividades Desempenhadas:}

\begin{itemize}
\item
  Participação em reunião com os AGES II para redefinir e redistribuir
  as tasks de testes.
\item
  Solução de conflitos de merge relacionados a alterações simultâneas em
  componentes e arquivos de teste.
\item
  Continuidade no estudo da estrutura de testes (src/test/msw x
  tests/unit), reforçando quando utilizar mocks e quando testar
  comportamento real.
\item
  Implementação dos testes do ExperienceCard, cobrindo cenários de
  renderização e interações básicas.
\item
  Início dos testes da Tela de Reservas, explorando a simulação de
  comportamento do usuário (renderização do filtro, lista de
  experiências, elementos principais da página).
\item
  Prática recorrente de execução, depuração e ajuste de testes,
  aprimorando a leitura de mensagens de erro e logs do Vitest.
\end{itemize}

\textbf{Ferramentas Utilizadas:}

\begin{itemize}
\item
  Vitest (execução, watch e depuração de testes).
\item
  React Testing Library (simulação de interação e queries na tela).
\item
  MSW (estudo e uso como infraestrutura de mocks quando necessário).
\item
  Node.js, Vite (ambiente de desenvolvimento e testes).
\item
  Git/GitHub (merge, resolução de conflitos e versionamento dos arquivos
  de teste).
\end{itemize}

\textbf{Contribuições no desenvolvimento:}

\begin{itemize}
\item
  Ampliação da cobertura de testes em componentes centrais do fluxo de
  reservas (ExperienceCard e Tela de Reservas).
\item
  Maior estabilidade e confiança no frontend, graças à validação
  automatizada de comportamento e renderização.
\item
  Consolidação do entendimento sobre o ecossistema de testes do projeto,
  facilitando a continuidade do trabalho de QA nas próximas sprints.
\item
  Evolução pessoal em testes automatizados de frontend, reduzindo a
  insegurança inicial e preparando o caminho para atuar com mais
  autonomia nessa área.
\end{itemize}

\subsection{3.4. Conclusão}\label{conclusuxe3o-1}

O desenvolvimento do Pró-Mata marcou um avanço significativo na minha
formação, não apenas pelo volume de tecnologias novas que explorei, mas
pela maneira como aprendi a conectar cada parte do sistema. Se nas
primeiras sprints eu transformava protótipos em telas navegáveis, nas
sprints seguintes aprendi a compreender a engenharia por trás do
produto: estado, rotas, validação, testes, arquitetura de componentes e
padrões visuais. Essa combinação elevou meu entendimento de
desenvolvimento web para um nível que eu não tinha alcançado antes.

No aspecto técnico, a AGES II ampliou meu repertório de forma prática.
Conheci ferramentas que até então eram completamente novas para mim ---
Zustand, Zod, TanStack Router, TanStack Query, Axios --- e cada uma
delas passou a fazer sentido à medida que eu entendia seu papel dentro
da aplicação. O Tailwind, que eu inicialmente julgava confuso, se tornou
uma das ferramentas que mais gostei de aprender, acelerando minha forma
de pensar interface, padronizando cores, sombras, espaçamentos e
mantendo o design system coeso. Também vivenciei minha primeira
experiência consistente com testes de frontend. Foi uma jornada difícil
e frustrante, especialmente na Sprint 4, quando percebi a complexidade
real por trás de mocks, queries, coverage e simulação de interação com a
interface. Embora eu tenha recorrido ao apoio da IA para concluir parte
dos testes, essa experiência deixou claro o quanto eu preciso evoluir
nessa área --- e essa consciência, por si só, já foi um aprendizado
importante.

Do ponto de vista de soft skills, houve um crescimento marcante. A
diferença entre minha postura no AGES I e no AGES II foi evidente:
participei mais, perdi o receio de fazer perguntas, deixei de me sentir
deslocado tecnicamente e me tornei mais confiante para propor ideias e
assumir tarefas desafiadoras. A experiência de criar fluxos, protótipos
e decisões visuais no Figma reforçou meu papel no time e me fez perceber
como posso contribuir não apenas com código, mas também com visão de
produto. Ao mesmo tempo, desenvolvi uma comunicação mais madura,
melhorei minhas descrições de PRs e passei a manter um processo de
trabalho mais organizado e previsível.

Em relação às disciplinas cursadas, percebi uma conexão real com
conteúdos acadêmicos. Arquitetura e Projeto de Software ajudaram a
entender o sistema como um todo, seus módulos, responsabilidades e
padrões. Sistemas Operacionais, por outro lado, conectou vários pontos
sobre como navegadores lidam com processos, abas e concorrência --- algo
que, mesmo indiretamente, influenciou minha compreensão sobre estados,
carregamento e estrutura de aplicações web. Foi um semestre em que
estudei mais do que nunca, muitas vezes indo além do necessário, movido
pela curiosidade de entender como cada peça funcionava por dentro.

Sobre o projeto em si, o Pró-Mata foi, sem dúvidas, o projeto mais
extenso e com o maior escopo que já participei até agora. Foi uma
experiência organizada, visualmente interessante e tecnicamente
desafiadora, com espaço para evoluir tanto no design quanto na lógica.
Trabalhar nesse sistema ampliou meu repertório como desenvolvedor e me
permitiu vivenciar, de forma concreta, a construção de um produto que
exige planejamento, consistência e integração entre várias etapas do
desenvolvimento.

Como autocrítica, acredito que poderia ter pedido mais ajuda em alguns
momentos. Por mergulhar profundamente em cada tecnologia, acabei
prolongando tarefas que poderiam ter avançado mais rápido. Por outro
lado, esse tempo investido me deu um aprendizado sólido, e hoje, com
essa bagagem, sinto que está muito mais concreto meu entendimento de
tudo que compõe um sistema.

Agora sei construir componentes complexos, estruturar fluxos de
navegação, padronizar um design system, escrever PRs melhores,
raciocinar sobre arquitetura de estados e, mesmo lutando com testes, já
compreendo o caminho para dominá-los.

De forma geral, o que fica da AGES II é a sensação de evolução genuína.
Evoluí tecnicamente, evoluí como membro de equipe e evoluí no
entendimento de como transformar um sistema do Figma até o código final.
Sigo para as próximas etapas com mais confiança, mais repertório e uma
clareza maior sobre meu papel como desenvolvedor. Ao mesmo tempo,
reconheço que ainda me sinto apreensivo quanto à transição para o AGES
III, especialmente por não dominar completamente áreas como
infraestrutura e arquitetura. Ainda assim, vejo isso não como um
obstáculo, mas como um próximo passo natural: conhecimentos que posso
alcançar com estudo, prática e dedicação. Estou animado para o que vem a
seguir e confiante de que essa trajetória continuará ampliando minhas
habilidades e minha maturidade profissional.

\section{4 - PROJETOS AGES III -
``FINUP''}\label{projetos-ages-iii---finup}

\subsection{4.1 Introdução}\label{introduuxe7uxe3o-2}

A dificuldade de entender, organizar e melhorar a própria vida
financeira é uma dor muito comum no Brasil. Muitos brasileiros chegam ao
fim do mês sem saber exatamente para onde o dinheiro foi --- não por
falta de interesse, mas porque nunca foram ensinados, de forma simples e
prática, a cuidar do próprio dinheiro. Diante desse cenário, surgiu a
proposta do FinUp: um aplicativo mobile que funciona como um "copiloto
financeiro do brasileiro", ajudando o usuário a controlar seus gastos,
tomar melhores decisões no dia a dia e aprender finanças de forma leve,
gamificada e sem julgamento.

A solução une organização financeira, inteligência artificial e educação
financeira em uma experiência acessível, jovem e prática. Entre suas
principais funcionalidades estão o controle de receitas e despesas, um
diagnóstico financeiro inicial, um indicador de saúde financeira (FinUp
Score), um assistente de IA para orientação no cotidiano, trilhas de
educação financeira gamificadas e um painel administrativo web.

Mais do que uma ferramenta de controle de gastos, o FinUp representa uma
iniciativa de impacto social, buscando democratizar o acesso à educação
financeira e empoderar seus usuários para que construam uma relação mais
saudável e consciente com o próprio dinheiro.

\subsection{3.2 Desenvolvimento do
Projeto.}\label{desenvolvimento-do-projeto.-1}

\textbf{Stakeholders do Projeto:}

\begin{itemize}
\item
  Vitor Oppermann
\item
  Diogo Ninov
\end{itemize}

\textbf{Período de Execução}: Segundo semestre de 2026.

\textbf{Professor Orientador:} Afonso Sales

Figura 3 - Foto do time

\begin{enumerate}
\def\labelenumi{\arabic{enumi}.}
\item
  \textbf{Repositório do código Fonte do Projeto}
\end{enumerate}

Wiki: \url{https://tools.ages.pucrs.br/2026-2/6lmnp/finup/finup-wiki}

Frontend Web:
\url{https://tools.ages.pucrs.br/2026-2/6lmnp/finup/finup-web}

Frontend Mobile:
\url{https://tools.ages.pucrs.br/2026-2/6lmnp/finup/finup-mobile}

Backend:
\url{https://tools.ages.pucrs.br/2026-2/6lmnp/finup/finup-backend}

\begin{enumerate}
\def\labelenumi{\arabic{enumi}.}
\setcounter{enumi}{1}
\item
  \textbf{Banco de Dados utilizado}
\end{enumerate}

\begin{quote}
Link para a wiki do banco:
\url{https://tools.ages.pucrs.br/2026-2/6lmnp/finup/finup-wiki/-/wikis/Banco\%20de\%20Dados}

O banco de dados do FinUp utiliza o modelo relacional, organizado em
três domínios: perfil do usuário, controle financeiro e educação
financeira gamificada . A escolha se justifica pela natureza estruturada
e interligada dos dados, garantindo integridade e consistência por meio
de chaves primárias e estrangeiras. As convenções adotadas incluem
nomenclatura em inglês com PascalCase e campos de auditoria, assegurando
padronização e facilitando a evolução do banco.
\end{quote}

\begin{enumerate}
\def\labelenumi{\arabic{enumi}.}
\setcounter{enumi}{2}
\item
  \textbf{Arquitetura utilizada}
\end{enumerate}

Link para wiki da arquitetura:
\url{https://tools.ages.pucrs.br/2026-2/6lmnp/finup/finup-wiki/-/wikis/arquitetura}

A arquitetura do FinUp foi dividida em quatro repositórios: finup-web
(React + Vite) para o painel administrativo, finup-mobile (React Native
+ Expo) para o aplicativo do usuário, finup-backend (Java Spring +
Python para IA) para a API e lógica de negócio, e finup-infra para
gerenciamento da infraestrutura {[}1{]}. A infraestrutura é hospedada na
AWS, utilizando EC2 para os containers Docker (frontend web, backend e
PostgreSQL), Cognito para autenticação de usuários, Bedrock como motor
central de IA (chatbot, trilhas e análise financeira), Lambda para
requisições específicas do painel admin, e EBS para armazenamento do
banco de dados. Além disso, foi adotado o Railforge como ferramenta de
padronização do fluxo de desenvolvimento nas diferentes stacks do
projeto.

\begin{enumerate}
\def\labelenumi{\arabic{enumi}.}
\setcounter{enumi}{3}
\item
  \textbf{Protótipos das telas desenvolvidas}
\end{enumerate}

https://tools.ages.pucrs.br/2026-2/6lmnp/finup/finup-wiki/-/wikis/Design(Mockups)

Os mockups do projeto foram feitos no Figma, exibindo o design e fluxo
das interfaces.

\begin{enumerate}
\def\labelenumi{\arabic{enumi}.}
\setcounter{enumi}{4}
\item
  \textbf{Tecnologias Utilizadas}
\end{enumerate}

Web (finup-web): React, Vite, Tailwind CSS, TanStack Router, Axios.

Mobile (finup-mobile): React Native, Expo, Expo Router.

\begin{quote}
Backend (finup-backend): Java, Spring Boot, Python, PostgreSQL,
Swagger/OpenAPI.

Infraestrutura e DevOps (finup-infra): Docker, Docker Compose, GitHub
Actions (CI/CD), Railforge.

Serviços AWS: EC2 (hospedagem dos containers), Cognito (autenticação),
Bedrock (motor de IA), Lambda (requisições do painel admin), EBS
(armazenamento do banco).
\end{quote}

\subsection{3.3 Atividades desempenhadas pelo aluno no
projeto}\label{atividades-desempenhadas-pelo-aluno-no-projeto-2}

\textbf{3.3.1 Sprint 0}

A Sprint 0 foi dedicada à definição das bases arquiteturais e
organizacionais do projeto FinUp, incluindo a análise da arquitetura do
sistema, a estruturação dos repositórios e ambientes de desenvolvimento,
e o estudo da infraestrutura em nuvem que sustentará a aplicação.

\textbf{Atividades Desempenhadas:}

\begin{itemize}
\item
  Análise e definição da arquitetura inicial do FinUp, considerando as
  camadas web, mobile, backend e infraestrutura.
\item
  Estudo e decisão sobre a estratégia de organização dos repositórios
  (monorepo x multirepo), levando em conta a divisão entre as equipes
  AGES II, III e IV.
\item
  Criação dos repositórios Web, Mobile, Backend e Infra no GitHub.
\item
  Configuração das branches main e develop com proteções, regras de
  approval e squash merge.
\item
  Organização dos times e definição de permissões de acesso por equipe
  no GitHub.
\item
  Estudo comparativo de serviços AWS para hospedagem do backend (EC2,
  ECS Fargate e Elastic Beanstalk).
\item
  Início da configuração do LocalStack para simulação e testes locais de
  serviços AWS.
\item
  Estudos complementares sobre Spring Boot, CI/CD e coleta de dados
  financeiros.
\end{itemize}

\textbf{Ferramentas Utilizadas:}

\begin{itemize}
\item
  GitHub (repositórios, branches, permissões), AWS (EC2, ECS Fargate,
  Elastic Beanstalk), LocalStack.
\end{itemize}

\textbf{Contribuições no desenvolvimento:}

Esta sprint representou o alicerce técnico e organizacional do projeto.
A principal contribuição foi estabelecer uma estrutura de repositórios
clara e bem protegida, garantindo que as três equipes pudessem trabalhar
de forma paralela e organizada. O estudo dos serviços AWS permitiu uma
visão inicial sobre as possibilidades de hospedagem e deploy, enquanto a
configuração do LocalStack abriu caminho para testes locais sem depender
de recursos em nuvem. Um desafio relevante foi a indefinição inicial
sobre a divisão das equipes, que impactou diretamente a escolha entre
monorepo e multirepo, reforçando o aprendizado de que decisões
arquiteturais estão intimamente ligadas à organização das pessoas no
projeto.

\textbf{3.3.2 Sprint 1}

A Sprint 1 foi dedicada à padronização do fluxo de desenvolvimento com a
adoção do Railforge, ao aprofundamento em Docker e configuração do
ambiente local do backend, à implantação das pipelines de CI nos
repositórios e ao apoio na integração dos primeiros fluxos funcionais do
aplicativo mobile.

\textbf{Atividades Desempenhadas:}

\begin{itemize}
\item
  Análise da adoção do Railforge no fluxo de desenvolvimento do FinUp,
  abrangendo requisitos, escopo, validações e commits.
\item
  Definição do Railforge em modo educativo, sem bloqueios rígidos para
  os desenvolvedores.
\item
  Estruturação documental e criação das variantes Basic, Backend,
  Frontend e Mobile do Railforge, com customização conforme cada stack.
\item
  Estudo e aplicação de conceitos de Docker, incluindo imagens,
  containers, portas, redes, Dockerfile e Docker Compose.
\item
  Configuração do ambiente local do backend com PostgreSQL e execução da
  API Spring Boot.
\item
  Validação dos endpoints de health, Swagger e OpenAPI.
\item
  Configuração da CI para ser disparada ao commitar em qualquer branch,
  com abertura de PRs nos repositórios Mobile, Web e Backend.
\item
  Análise da navegação do FinUp Mobile com Expo Router.
\item
  Integração e validação dos fluxos de login e cadastro no mobile.
\item
  Resolução de conflitos com a branch develop, ajustes de dependências e
  correção da dependência expo-font no CI.
\item
  Realização de code reviews em PRs da equipe.
\end{itemize}

\textbf{Ferramentas Utilizadas:}

GitHub (PRs, branches, CI), Docker, Docker Compose, Railforge, Spring
Boot, PostgreSQL, Swagger/OpenAPI, Expo Router.

\textbf{Contribuições no desenvolvimento:}

Esta sprint marcou a transição do planejamento para a execução prática
do projeto. A principal contribuição foi a adoção e customização do
Railforge, criando variantes específicas para cada stack e garantindo
consistência no fluxo de trabalho das equipes. A configuração do
ambiente local com Docker e a validação da API Spring Boot permitiram
que o backend fosse executado de forma independente por qualquer
desenvolvedor. A implantação das pipelines de CI trouxe automação e
qualidade, embora tenha exigido a resolução de problemas como falhas com
expo-font e conflitos no package-lock. No mobile, o apoio na integração
dos fluxos de login e cadastro contribuiu para a estabilização das
primeiras telas funcionais. Entre os principais aprendizados, destaca-se
que ferramentas de padronização precisam ser adaptadas ao contexto de
cada stack, e que o domínio de Docker e CI são fundamentais para a
produtividade em equipes distribuídas.

\section{REFERÊNCIAS}\label{referuxeancias}

\begin{itemize}
\item
  React: \url{https://react.dev}
\item
  TypeScript: \url{https://www.typescriptlang.org/docs/}
\item
  Tailwind CSS: \url{https://tailwindcss.com/docs}
\item
  TanStack Router: \url{https://tanstack.com/router/latest}
\item
  Zustand:
  \url{https://docs.pmnd.rs/zustand/getting-started/introduction}
\item
  Zod: \url{https://zod.dev}
\item
  Vite: \url{https://vitejs.dev/guide/}
\item
  Git --- Documentação: \url{https://git-scm.com/doc}
\item
  React Native: \url{https://reactnative.dev/docs/getting-started}
\item
  Expo: \url{https://docs.expo.dev}
\item
  Expo Router: \url{https://docs.expo.dev/router/introduction/}
\item
  Axios: \url{https://axios-http.com/docs/intro}
\item
  Java: \url{https://docs.oracle.com/en/java/}
\item
  Spring Boot:
  \url{https://docs.spring.io/spring-boot/documentation.html}
\item
  \href{PostgreSQL:\%20https://www.postgresql.org/docs/}{PostgreSQL:
  https://www.postgresql.org/docs/}
\item
  Swagger / OpenAPI: https://swagger.io/docs/
\item
  Docker: \url{https://docs.docker.com}
\item
  AWS EC2: \url{https://docs.aws.amazon.com/ec2/}
\item
  AWS Cognito: \url{https://docs.aws.amazon.com/cognito/}
\item
  AWS Bedrock: \url{https://docs.aws.amazon.com/bedrock/}
\item
  AWS Lambda: \url{https://docs.aws.amazon.com/lambda/}
\item
  AWS EBS: \url{https://docs.aws.amazon.com/ebs/}
\item
  LocalStack: \url{https://docs.localstack.cloud/overview/}
\end{itemize}

\end{document}
